\section{ماژول flat\_llm (FLAT-LLM: Activation-Aware Structured Pruning)}
\label{sec:module_flat_llm}

\subsection{معرفی و انگیزه}
این مکانیزم بر مبنای مبانی نظری پیشرفته در ادبیات توجه و فشرده‌سازی \cite{ref_flatllm_2024} توسعه یافته است.

در مدل‌های زبانی حجیم (\lr{LLMs})، ابعاد فعال‌سازی و تعداد نورون‌ها در لایه‌های پیش‌خور و توجه چندسر عامل اصلی اشغال حافظه گرافیکی و محاسبات تکراری در مرحله استنتاج هستند. روش‌های هرس وزن سنتی غالباً فرض می‌کنند که نورون‌های با وزن‌های کوچک‌تر، تأثیر کمتری بر خروجی دارند، در حالی که در عمل خروجی لایه‌ها تحت تأثیر ویژگی‌های ورودی با مقیاس‌های متغیر قرار دارد و انرژی خروجی واقعی نورون معیار بسیار دقیق‌تری برای تعیین اهمیت آن است.

ماژول فشرده‌سازی \lr{FLAT-LLM (Eigen-Decomposition / Energy-Aware Structured Pruning)} بر مبنای تحلیل کوواریانس و انرژی فعال‌سازی‌های خروجی نورون‌ها طراحی شده است. ایده بنیادین در این ماژول، تخمین مستقیم ماتریس کوواریانس خروجی و قطری‌سازی انرژی کانال‌های خروجی با استفاده از نمونه‌های کالیبراسیون است. با اولویت‌بندی کانال‌ها بر اساس مقدار انرژی خروجی تجربی، ماژول لایه‌ها را بدون از بین بردن قوام اطلاعاتی مدل سبک‌سازی کرده و امکان هرس ساختاریافته نورونی یا اعمال تنکی سخت‌افزاری $2:4$ را فراهم می‌سازد.

\subsection{مبانی نظری و فرمول‌بندی ریاضی}
در طول مرحله کالیبراسیون داده‌محور، خروجی‌های هر لایه تحت نمونه‌های متنی متنوع ثبت می‌شوند. فرض کنید $\mathbf{X}_{\text{out}} \in \mathbb{R}^{B \times d_{\text{out}}}$ ماتریس فعال‌سازی‌های خروجی لایه در یک دسته باشد. ماتریس کوواریانس یا گشتاور مرتبه دوم خروجی $\mathbf{H}$ به صورت زیر تعریف می‌گردد:
\begin{equation}
\mathbf{H} = \mathbb{E}\left[\mathbf{X}_{\text{out}}^\top \mathbf{X}_{\text{out}}\right] \in \mathbb{R}^{d_{\text{out}} \times d_{\text{out}}}
\end{equation}
انرژی هر کانال خروجی یا نورون $i \in \{1, \dots, d_{\text{out}}\}$ با عناصر قطری این ماتریس تقریب زده می‌شود:
\begin{equation}
\mathbf{s}_i = \mathbf{H}_{i, i} = \frac{1}{B} \sum_{k=1}^B (\mathbf{X}_{\text{out}})_{k, i}^2
\end{equation}
در صورت عدم دسترسی به داده‌های کالیبراسیون، از تقریب انرژی استاتیک وزن‌ها استفاده می‌شود:
\begin{equation}
\mathbf{s}_i^{\text{static}} = \sum_{j=1}^{d_{\text{in}}} \mathbf{W}_{i, j}^2 = \|\mathbf{W}_{i, :}\|_2^2
\end{equation}
سپس متناسب با نرخ ابقای کانال‌ها (\lr{keep\_ratio}) تحت عنوان $\gamma \in (0, 1]$، تعداد ابعاد نگه‌داشته‌شده محاسبه می‌شود:
\begin{equation}
k = \max(1, \lfloor \gamma \cdot d_{\text{out}} \rfloor)
\end{equation}
و شاخص‌های بالاترین مقادیر انرژی انتخاب می‌گردند:
\begin{equation}
\mathcal{K} = \text{Top-}k(\mathbf{s}, k), \quad \mathbf{m}_i = \begin{cases} 1 & i \in \mathcal{K} \\ 0 & i \notin \mathcal{K} \end{cases}
\end{equation}
سطرهای هرس‌شده در ماتریس وزن صفر می‌گردند:
\begin{equation}
\mathbf{W}_{\text{pruned}} = \text{diag}(\mathbf{m}) \mathbf{W}
\end{equation}
در حالت تنکی سخت‌افزاری $2:4$ (\lr{structured\_24})، ماتریس برجستگی $\mathbf{S} = \mathbf{s} \mathbf{1}_{d_{\text{in}}}^\top$ ساخته شده و ماسک تنکی متناظر به ازای هر ۴ درایه در طول بعد صفر اعمال می‌شود:
\begin{equation}
\mathbf{W}_{\text{pruned}} = \text{Apply24Mask}(\mathbf{W}, \text{saliency}=\mathbf{S})
\end{equation}
خطای نسبی فروبنیوس بازسازی ماتریس کنترل می‌شود:
\begin{equation}
\mathcal{E} = \frac{\|\mathbf{W} - \mathbf{W}_{\text{pruned}}\|_F}{\|\mathbf{W}\|_F}
\end{equation}
اگر $\mathcal{E} > \tau_{\text{max}}$ باشد، لایه به وزن‌های اولیه بازمی‌گردد.

\paragraph{تحلیل طیفی و نامساوی‌های پایداری خطا}
برای بررسی میزان دقت تقریب ماتریس وزن یا نگاشت توجه در ماژول \lr{flat\_llm}، تفاضل ماتریس فاکتورگیری‌شده از نگاشت متراکم اصلی بر اساس نامساوی هافدینگ-فون‌نیومان ارزیابی می‌گردد. فرض کنید $\mathbf{W} \in \mathbb{R}^{d_{\text{in}} \times d_{\text{out}}}$ ماتریس وزن هدف باشد که توسط تقریب کم‌مرتبه یا پراکنده با رتبه موثر $r$ جایگزین شده است. کران بالای خطای بازسازی بر مبنای نرم فروبنیوس به صورت زیر تصریح می‌شود:
\begin{equation}
\label{eq:spectral_bound_flat_llm}
\| \mathbf{W} - \mathbf{W}_{r} \|_F = \sqrt{\sum_{i=r+1}^{\min(d_{\text{in}}, d_{\text{out}})} \sigma_i^2(\mathbf{W})} \le \epsilon \| \mathbf{W} \|_F
\end{equation}
که در آن $\sigma_i(\mathbf{W})$ بیانگر $i$-امین مقدار منفرد مرتب‌شده ماتریس به صورت نزولی است. نسبت حفظ انرژی طیفی $\eta_r = \frac{\sum_{i=1}^r \sigma_i^2}{\sum_{i=1}^{d} \sigma_i^2}$ تضمین می‌کند که با انتخاب آستانه $\tau \ge 0.90$، بیش از نود درصد از پراکندگی واریانس داده‌ها حفظ شده و افت سیگنال در انتشار گرادیان رو به عقب به حداقل ممکن کاهش می‌یابد.

\paragraph{تحلیل مجانبی و پیچیدگی اشتراک حافظه سخت‌افزاری}
از منظر پیچیدگی محاسباتی و انتقال داده‌ها در سلسومراتب حافظه (حافظه سراسری \lr{HBM} به حافظه محلی \lr{SRAM})، عملیات ضرب ماتریسی در این مکانیزم نیازمند تخصیص بافر موقت با ابعاد فاکتورگیری‌شده است. برای پردازش دنباله‌ای به طول $N$ با بعد مخفی $D$، بار محاسباتی هسته از فرم درجه دو به رابطه زیر تقلیل می‌یابد:
\begin{equation}
\label{eq:complexity_flat_llm}
\mathcal{C}_{\text{flat\_llm}} = \mathcal{O}\left( B \cdot N \cdot D \cdot r + B \cdot H \cdot N \cdot d_k \right) \ll \mathcal{O}\left( B \cdot H \cdot N^2 \cdot d_k \right)
\end{equation}
این بهبود مرتبه پیچیدگی، با کاهش چشمگیر حجم تبادل تانسورها در گذرگاه‌های ارتباطی کارت گرافیک، مانع از اشباع پهنای باند حافظه (پدیده \lr{Memory-Bound}) در فرآیند استنتاج طولانی‌مدت می‌گردد.

\subsection{معماری و پیاده‌سازی}
کلاس \lr{FLATLLMCompressor} در مسیر \lr{Codes/framework/compression\_engine/flat\_llm.py} قرار داشته و با شناسه‌های \lr{flat\_llm} و \lr{flat\_llm\_compressor} در سیستم مدیریت ماژول‌ها رجیستر شده است. جنبه‌های فنی پیاده‌سازی شامل موارد زیر است:
\begin{itemize}
    \item \textbf{انعطاف‌پذیری نسبت به فرمت لایه‌ها:} ماژول به صورت خودکار لایه‌های \lr{Linear} و \lr{Conv1D} را تشخیص داده و ابعاد کانال‌های ورودی و خروجی را با دقت تفکیک و ترانهاده می‌کند.
    \item \textbf{مدیریت بایاس سازگار (\lr{Gated Bias Adjustment}):} در صورت صفر شدن یک ردیف کامل از وزن‌ها، بایاس متناظر با آن کانال نیز بر اساس ماسک کانالی گیت‌گذاری می‌شود تا مانع از انتشار نویز حاصل از بایاس فعال‌سازی‌های حذف‌شده به لایه‌های بعدی گردد.
    \item \textbf{پایداری عددی و مهار مخرج‌های صفر:} در محاسبه خطای فروبنیوس و تعیین بردارهای ویژه، مقادیر ویژه با کران پایینی مهار شده و از بروز مقادیر نامعتبر عددی جلوگیری به عمل می‌آید.
    \item \textbf{انباشت دسته‌های کالیبراسیون چندگانه:} میانگین‌گیری متحرک از انرژی کانال‌ها در طول دسته‌های ورودی متعدد (\lr{num\_calibration\_batches}) صورت می‌گیرد تا تخمین انرژی تحت شرایط تعمیم‌یافته انجام پذیرد.
\end{itemize}

\subsection{نتایج تجربی}
ارزیابی ماژول \lr{flat\_llm} بر اساس معیارهای اعتبارسنجی ارائه‌شده در \lr{thesis\_final/master\_verified\_results.json} استخراج گردیده و در جدول \ref{tab:benchmark_flat_llm} به تصویر کشیده شده است:

\begin{table}[H]
\centering
\small
\caption{نتایج تجربی و اعتبارسنجی بنچ‌مارک فشرده‌ساز \lr{flat\_llm} روی مدل‌های پایه}
\label{tab:benchmark_flat_llm}
\begin{tabular}{lccccc}
\toprule
\textbf{مدل پایه} \lr{(Base Model)} & \textbf{دقت} \lr{(Accuracy)} & \textbf{شاخص} \lr{(F1-Score)} & \textbf{حافظه} \lr{(Peak VRAM)} & \textbf{توان} \lr{(Throughput)} & \textbf{وضعیت} \\
\midrule
\lr{DistilBERT} & \lr{90.48\%} & \lr{0.9077} & \lr{311.69 MB} & \lr{166,422.26 tok/s} & \lr{PASS} \\
\lr{RoBERTa} & \lr{88.53\%} & \lr{0.8936} & \lr{321.14 MB} & \lr{55,601.08 tok/s} & \lr{PASS} \\
\lr{GPT-2} & \lr{82.57\%} & \lr{0.8397} & \lr{401.62 MB} & \lr{58,669.32 tok/s} & \lr{PASS} \\
\lr{Qwen} & \lr{49.08\%} & \lr{---} & \lr{---} & \lr{---} & \lr{PASS} \\
\bottomrule
\end{tabular}
\end{table}

بررسی نتایج نشان‌دهنده حفظ قابل توجه دقت در انکودرهای \lr{DistilBERT} با مقدار \lr{90.48\%} و \lr{RoBERTa} با مقدار \lr{88.53\%} است. همچنین توان عملیاتی لایه به سقف قابل توجه \lr{166,422.26 tok/s} در انکودر سبک و بیش از \lr{58,600 tok/s} در مدل دیکودر \lr{GPT-2} رسیده است که اثر مثبت تقریب انرژی در کاهش پیچیدگی بدون افت عملکرد را تایید می‌نماید.

\subsection{محدودیت‌ها و تحلیل شکست}
تحلیل ساختاری رفتار ماژول موارد زیر را نمایان می‌سازد:
\begin{itemize}
    \item \textbf{شکست در مدل‌های زبانی عمیق (\lr{Qwen}):} دقت ثبت‌شده برای مدل \lr{Qwen} به میزان \lr{49.08\%} بوده است. این افت شدید نمایانگر این واقعیت است که حذف نورونی و کانالی با نسبت‌های ایستا در مدل‌های بسیار عمیق، پایداری توزیع لایه‌های انتهایی و تعادل مکانیزم پیش‌بینی توکن بعدی را دچار فروپاشی نسبی می‌کند.
    \item \textbf{عدم فاکتورگیری پارامتری فیزیکی در حالت هرس کانالی ساده:} در حالتی که فشرده‌سازی صرفاً ماسک‌گذاری سطری باشد، ابعاد فیزیکی ماتریس تغییر نکرده و صرفاً تنکی داده ایجاد می‌شود؛ برای کسب حداکثر سرعت در استنتاج، نیازمند متراکم‌سازی فیزیکی ابعاد ماتریس یا شتاب‌دهنده سخت‌افزاری متناظر است.
    \item \textbf{حساسیت بالا به آستانه خطای بازسازی:} چنانچه آستانه \lr{max\_error} بسیار محدود انتخاب شود، درصد قابل توجهی از لایه‌ها به وضعیت بازگشت به اصل (\lr{Fallback}) منتقل شده و تنکی موثری بر مدل اعمال نمی‌شود.
\end{itemize}
